\documentclass[10pt,a4paper]{article}

% -------------------------------------------------- %
% Packages                                           %
% -------------------------------------------------- %
\usepackage{fontspec}
[% if taal == 'en' %]
\usepackage[english]{babel}
[% else %]
\usepackage[dutch]{babel}
[% endif %]
\usepackage{geometry}
\geometry{a4paper, margin=2.5cm}
\usepackage{graphicx}
\usepackage{xcolor}
\usepackage{hyperref}
\usepackage{fancyhdr}
\usepackage{array}
\usepackage{longtable}
\usepackage{booktabs}
\usepackage{ifthen}

% Rapport package (bevindingen met CVSS-kleuren)
\usepackage{opmaak/packages/rapport}

% -------------------------------------------------- %
% Metadata                                           %
% -------------------------------------------------- %
\newcommand{\rapporttitel}{[[ titel ]]}
\newcommand{\rapportsubtitel}{[[ subtitel ]]}
\newcommand{\rapportauteur}{[[ auteur ]]}
\newcommand{\rapportproject}{[[ project ]]}
\newcommand{\rapportdatum}{[[ datum ]]}
\newcommand{\rapportclassificatie}{[[ classificatie ]]}
\newcommand{\rapporttlpkleur}{[[ tlp_color ]]}

\hypersetup{
  pdftitle={\rapporttitel},
  pdfauthor={\rapportauteur},
}

% -------------------------------------------------- %
% Pagina-opmaak                                      %
% -------------------------------------------------- %
\pagestyle{fancy}
\fancyhf{}
\fancyhead[L]{\small \rapporttitel}
\fancyhead[R]{\small \textcolor{\rapporttlpkleur}{\rapportclassificatie}}
\fancyfoot[L]{\small \rapportauteur}
\fancyfoot[C]{\small \rapportproject}
\fancyfoot[R]{\small \thepage}
\renewcommand{\headrulewidth}{0.4pt}
\renewcommand{\footrulewidth}{0.4pt}

\graphicspath{[[ '{' ]]evidence/[[ '}' ]]}

% -------------------------------------------------- %
% Document                                           %
% -------------------------------------------------- %
\begin{document}

% --- Voorblad ---
\begin{titlepage}
  \centering
  \vspace*{4cm}
  {\Huge\bfseries \rapporttitel \par}
  \vspace{1cm}
  {\Large \rapportsubtitel \par}
  \vspace{2cm}
  {\large \rapportproject \par}
  \vspace{1cm}
  {\large \rapportauteur \par}
  \vspace{1cm}
  {\large \rapportdatum \par}
  \vspace{2cm}
  [% if is_draft %]
  [% if taal == 'en' %]
  {\Large\textcolor{ored}{DRAFT} \par}
  [% else %]
  {\Large\textcolor{ored}{CONCEPT} \par}
  [% endif %]
  [% endif %]
  \vfill
  [% if taal == 'en' %]
  {\small Classification: \textcolor{\rapporttlpkleur}{\rapportclassificatie} \par}
  [% else %]
  {\small Classificatie: \textcolor{\rapporttlpkleur}{\rapportclassificatie} \par}
  [% endif %]
\end{titlepage}

% --- Inhoudsopgave ---
\tableofcontents
\newpage

% --- Inleiding ---
[% if taal == 'en' %]
\section{Introduction}
[% else %]
\section{Inleiding}
[% endif %]

[% if testtype == 'redteam' %]
[% if taal == 'en' %]
This report contains the findings of the red team assessment performed on \rapportproject.
The objective was to simulate a realistic attack to evaluate the detection and response capabilities
of the organisation, as well as to identify vulnerabilities in people, processes and technology.
[% else %]
Dit rapport bevat de bevindingen van het red team assessment uitgevoerd op \rapportproject.
Het doel was het simuleren van een realistische aanval om de detectie- en responscapaciteiten
van de organisatie te evalueren, evenals het identificeren van kwetsbaarheden in mensen,
processen en technologie.
[% endif %]
[% else %]
[% if taal == 'en' %]
This report contains the findings of the penetration test performed on \rapportproject.
The objective was to identify vulnerabilities that could compromise the confidentiality,
integrity and availability of the systems.
[% else %]
Dit rapport bevat de bevindingen van de penetratietest uitgevoerd op \rapportproject.
Het doel van de test was het identificeren van kwetsbaarheden die de vertrouwelijkheid,
integriteit en beschikbaarheid van de systemen kunnen aantasten.
[% endif %]
[% endif %]

[% if taal == 'en' %]
The test was performed by \rapportauteur{} on \rapportdatum.
[% else %]
De test is uitgevoerd door \rapportauteur{} op \rapportdatum.
[% endif %]

\subsection{Scope}

[% if taal == 'en' %]
The following systems were part of the test:
[% else %]
De volgende systemen zijn onderdeel van de test:
[% endif %]
\begin{itemize}
[% for loc in scoped_hosts %]
  \item [[ loc | latex_escape ]]
[% endfor %]
[% if not scoped_hosts %]
[% if taal == 'en' %]
  \item See findings for specific locations
[% else %]
  \item Zie bevindingen voor specifieke locaties
[% endif %]
[% endif %]
\end{itemize}

[% if management_samenvatting %]
% --- Managementsamenvatting ---
[% if taal == 'en' %]
\section{Executive Summary}
[% else %]
\section{Managementsamenvatting}
[% endif %]

[[ management_samenvatting ]]

[% endif %]
% --- Methodologie ---
[% if taal == 'en' %]
\section{Methodology}
[% else %]
\section{Methodologie}
[% endif %]

[% if testtype == 'redteam' %]
[% if taal == 'en' %]
The red team assessment was performed in accordance with the MITRE ATT\&CK framework and the Cyber Kill Chain.
The test consisted of the following phases:

\begin{enumerate}
  \item \textbf{Reconnaissance} -- OSINT and target identification
  \item \textbf{Weaponisation} -- Prepare tooling and payloads
  \item \textbf{Delivery} -- Obtain initial access
  \item \textbf{Exploitation} -- Exploit vulnerabilities
  \item \textbf{Installation} -- Establish persistence
  \item \textbf{Command \& Control} -- Set up C2 channel
  \item \textbf{Actions on Objectives} -- Achieve goals
\end{enumerate}
[% else %]
Het red team assessment is uitgevoerd conform het MITRE ATT\&CK framework en de Cyber Kill Chain.
De test bestond uit de volgende fasen:

\begin{enumerate}
  \item \textbf{Verkenning} -- OSINT en doelidentificatie
  \item \textbf{Bewapening} -- Tooling en payloads voorbereiden
  \item \textbf{Aflevering} -- Initiele toegang verkrijgen
  \item \textbf{Exploitatie} -- Kwetsbaarheden benutten
  \item \textbf{Installatie} -- Persistentie vestigen
  \item \textbf{Commando \& Controle} -- C2-kanaal opzetten
  \item \textbf{Acties op doel} -- Doelstellingen realiseren
\end{enumerate}
[% endif %]
[% else %]
[% if taal == 'en' %]
The penetration test was performed in accordance with the OWASP Testing Guide and PTES (Penetration Testing
Execution Standard). The test consisted of the following phases:

\begin{enumerate}
  \item \textbf{Reconnaissance} -- Gather information about the target environment
  \item \textbf{Scanning} -- Identify services and vulnerabilities
  \item \textbf{Exploitation} -- Confirm vulnerabilities
  \item \textbf{Post-exploitation} -- Determine impact
  \item \textbf{Reporting} -- Document findings
\end{enumerate}
[% else %]
De penetratietest is uitgevoerd conform de OWASP Testing Guide en PTES (Penetration Testing
Execution Standard). De test bestond uit de volgende fasen:

\begin{enumerate}
  \item \textbf{Verkenning} -- Informatie verzamelen over de doelomgeving
  \item \textbf{Scanning} -- Diensten en kwetsbaarheden identificeren
  \item \textbf{Exploitatie} -- Kwetsbaarheden bevestigen
  \item \textbf{Post-exploitatie} -- Impact bepalen
  \item \textbf{Rapportage} -- Documenteren
\end{enumerate}
[% endif %]
[% endif %]

[% for scope in testscopes %]
[% if scope == 'whitebox' %]
[% if taal == 'en' %]
The test was (partially) performed as a white box assessment. The testers had full access to source code,
architecture documentation and configurations. This allows for a thorough analysis of all security layers.
[% else %]
De test is (deels) uitgevoerd als white box assessment. De testers hadden volledige toegang tot broncode,
architectuurdocumentatie en configuraties. Dit maakt een grondige analyse mogelijk van alle
beveiligingslagen.
[% endif %]
[% elif scope == 'greybox' %]
[% if taal == 'en' %]
The test was (partially) performed as a grey box assessment. The testers had limited information, such as
network topology, user accounts or API documentation. This simulates an attacker with limited internal knowledge.
[% else %]
De test is (deels) uitgevoerd als grey box assessment. De testers hadden beperkte informatie, zoals
netwerktopologie, gebruikersaccounts of API-documentatie. Dit simuleert een aanvaller met
beperkte interne kennis.
[% endif %]
[% else %]
[% if taal == 'en' %]
The test was (partially) performed as a black box assessment. The testers had no prior information about the
internal workings, source code or architecture of the target environment. This simulates the perspective of
an external attacker.
[% else %]
De test is (deels) uitgevoerd als black box assessment. De testers hadden vooraf geen informatie over de
interne werking, broncode of architectuur van de doelomgeving. Dit simuleert het perspectief van
een externe aanvaller.
[% endif %]
[% endif %]
[% endfor %]

% --- Tools en technieken ---
[% if taal == 'en' %]
\section{Tools and Techniques}

For this assessment, use was made of \textbf{Incompetent Bastard}, an open-source project
that provides support for security assessments such as the one described in this report.
In addition, the following tools were used during the assessment:

\begin{itemize}
  \item \textbf{Nmap} -- Network discovery and port scanning
  \item \textbf{Burp Suite} -- Web application security testing
  \item \textbf{Metasploit Framework} -- Exploit development and execution
  \item \textbf{BloodHound} -- Active Directory attack path analysis
  \item \textbf{CrackMapExec} -- Network penetration testing
  \item \textbf{Impacket} -- Network protocol interaction
\end{itemize}

The specific tools used may vary depending on the scope and nature of the assessment.
All tools are used in accordance with the agreed Rules of Engagement.
[% else %]
\section{Tools en technieken}

Voor dit onderzoek is gebruik gemaakt van \textbf{Incompetent Bastard}, een opensourceproject
dat ondersteuning biedt bij beveiligingsonderzoeken zoals uitgewerkt in dit rapport.
Daarnaast zijn de volgende tools ingezet tijdens het onderzoek:

\begin{itemize}
  \item \textbf{Nmap} -- Netwerk discovery en poortscan
  \item \textbf{Burp Suite} -- Webapplicatie beveiligingstesten
  \item \textbf{Metasploit Framework} -- Exploit ontwikkeling en uitvoering
  \item \textbf{BloodHound} -- Active Directory aanvalspad-analyse
  \item \textbf{CrackMapExec} -- Netwerk penetratietesten
  \item \textbf{Impacket} -- Netwerkprotocol interactie
\end{itemize}

De specifieke tools die worden ingezet kunnen varieren afhankelijk van de scope en aard van het onderzoek.
Alle tools worden ingezet conform de overeengekomen Rules of Engagement.
[% endif %]

[% if testtype == 'redteam' and (legsup_startpunt or legsup_privileges or legsup_informatie) %]
% --- Legs Up ---
\section{Legs Up}

[% if legsup_startpunt %]
[% if taal == 'en' %]
\subsection{Assumed Breach Starting Points}
[% else %]
\subsection{Assumed Breach Startpunten}
[% endif %]

[[ legsup_startpunt | latex_escape ]]

[% endif %]
[% if legsup_privileges %]
[% if taal == 'en' %]
\subsection{Privilege Levels}
[% else %]
\subsection{Privilege Niveaus}
[% endif %]

[[ legsup_privileges | latex_escape ]]

[% endif %]
[% if legsup_informatie %]
[% if taal == 'en' %]
\subsection{Provided Information}
[% else %]
\subsection{Verstrekte Informatie}
[% endif %]

[[ legsup_informatie | latex_escape ]]

[% endif %]
[% endif %]

[% if testtype == 'redteam' and scenarios %]
% --- Scenario's ---
[% if taal == 'en' %]
\section{Scenarios}
[% else %]
\section{Scenario's}
[% endif %]

[% for sc in scenarios %]
\subsection{[[ sc.naam | latex_escape ]]}

[[ sc.beschrijving | latex_escape ]]

[% endfor %]
[% endif %]

% --- Risicoinschatting ---
[% if taal == 'en' %]
\section{Risk Assessment}

The severity of findings is determined using the CVSS 4.0 framework
(Common Vulnerability Scoring System). The following classification is used:

\begin{center}
\begin{tabular}{|l|c|l|}
  \hline
  \textbf{Classification} & \textbf{CVSS Score} & \textbf{Colour} \\
  \hline
  Critical & 9.0 -- 10.0 & \textcolor{odarkred}{Dark red} \\
  High & 7.0 -- 8.9 & \textcolor{ored}{Red} \\
  Medium & 4.0 -- 6.9 & \textcolor{oorange}{Orange} \\
  Low & 0.1 -- 3.9 & \textcolor{oyellow}{Yellow} \\
  Informational & 0.0 & \textcolor{ogreen}{Green} \\
  \hline
\end{tabular}
\end{center}
[% else %]
\section{Risicoinschatting}

De ernst van bevindingen wordt bepaald aan de hand van het CVSS 4.0 framework
(Common Vulnerability Scoring System). De volgende classificatie wordt gehanteerd:

\begin{center}
\begin{tabular}{|l|c|l|}
  \hline
  \textbf{Classificatie} & \textbf{CVSS Score} & \textbf{Kleur} \\
  \hline
  Kritiek & 9.0 -- 10.0 & \textcolor{odarkred}{Donkerrood} \\
  Hoog & 7.0 -- 8.9 & \textcolor{ored}{Rood} \\
  Midden & 4.0 -- 6.9 & \textcolor{oorange}{Oranje} \\
  Laag & 0.1 -- 3.9 & \textcolor{oyellow}{Geel} \\
  Nihil & 0.0 & \textcolor{ogreen}{Groen} \\
  \hline
\end{tabular}
\end{center}
[% endif %]

% --- Samenvatting bevindingen ---
[% if taal == 'en' %]
\section{Summary of Findings}

A total of \textbf{[[ findings|length ]]} findings were identified:

\begin{center}
\begin{tabular}{|l|c|}
  \hline
  \textbf{Classification} & \textbf{Count} \\
  \hline
  \textcolor{odarkred}{Critical} & [[ counts.critical ]] \\
  \textcolor{ored}{High} & [[ counts.high ]] \\
  \textcolor{oorange}{Medium} & [[ counts.medium ]] \\
  \textcolor{oyellow}{Low} & [[ counts.low ]] \\
  \textcolor{ogreen}{Informational} & [[ counts.none ]] \\
  \hline
  \textbf{Total} & \textbf{[[ findings|length ]]} \\
  \hline
\end{tabular}
\end{center}
[% else %]
\section{Samenvatting bevindingen}

In totaal zijn \textbf{[[ findings|length ]]} bevindingen geidentificeerd:

\begin{center}
\begin{tabular}{|l|c|}
  \hline
  \textbf{Classificatie} & \textbf{Aantal} \\
  \hline
  \textcolor{odarkred}{Kritiek} & [[ counts.critical ]] \\
  \textcolor{ored}{Hoog} & [[ counts.high ]] \\
  \textcolor{oorange}{Midden} & [[ counts.medium ]] \\
  \textcolor{oyellow}{Laag} & [[ counts.low ]] \\
  \textcolor{ogreen}{Nihil} & [[ counts.none ]] \\
  \hline
  \textbf{Totaal} & \textbf{[[ findings|length ]]} \\
  \hline
\end{tabular}
\end{center}
[% endif %]

\newpage

% --- Verkenning & ontdekking (notities) ---
[% if notes %]
[% if taal == 'en' %]
\section{Reconnaissance \& Discovery}
[% else %]
\section{Verkenning \& ontdekking}
[% endif %]

[% for note in notes %]
\subsection{[[ note.naam | latex_escape ]]}

[[ note.uitwerken | latex_escape ]]

[% endfor %]
\newpage
[% endif %]

% --- Bevindingen ---
[% if taal == 'en' %]
\section{Findings}
[% else %]
\section{Bevindingen}
[% endif %]

[% for severity_label, group in grouped_findings %]
[% if group %]
[% if taal == 'en' %]
\subsection{[[ severity_label ]] findings}
[% else %]
\subsection{[[ severity_label ]] bevindingen}
[% endif %]

[% for item in group %]
\ibfinding{[[ item.id ]]}%
  {[[ item.titel | latex_escape ]]}%
  {[[ item.basescore ]]}%
  {[[ item.beschrijving | latex_escape ]]}%
  {[[ item.impact | latex_escape ]]}%
  {[[ item.aanbeveling | latex_escape ]]}

[% if item.locatie %]
[% if taal == 'en' %]
\noindent\textbf{Location:} \texttt{[[ item.locatie | latex_escape ]]}
[% else %]
\noindent\textbf{Locatie:} \texttt{[[ item.locatie | latex_escape ]]}
[% endif %]
\vspace{0.3em}
[% endif %]

[% if item.owasp %]
\noindent\textbf{OWASP:} [[ item.owasp | latex_escape ]] \quad
[% endif %]
[% if item.cwe %]
\textbf{CWE:} [[ item.cwe | latex_escape ]] \quad
[% endif %]
[% if item.cvss_vector %]
\textbf{CVSS Vector:} \texttt{\small [[ item.cvss_vector | latex_escape ]]}
[% endif %]

[% for ev in item.evidence %]
\ibevidence{[[ ev.path ]]}{[[ ev.caption | latex_escape ]]}
[% endfor %]

\newpage
[% endfor %]
[% endif %]
[% endfor %]

% --- Aanbevelingen ---
[% if taal == 'en' %]
\section{Recommendations Overview}

\begin{longtable}{|c|p{5cm}|c|p{7cm}|}
  \hline
  \textbf{ID} & \textbf{Finding} & \textbf{CVSS} & \textbf{Recommendation} \\
  \hline
  \endhead
[% else %]
\section{Overzicht aanbevelingen}

\begin{longtable}{|c|p{5cm}|c|p{7cm}|}
  \hline
  \textbf{ID} & \textbf{Bevinding} & \textbf{CVSS} & \textbf{Aanbeveling} \\
  \hline
  \endhead
[% endif %]
[% for item in findings %]
  [[ '%03d' | format(item.id) ]] & [[ item.titel | latex_escape ]] & [[ item.basescore ]] & [[ item.aanbeveling | latex_escape ]] \\
  \hline
[% endfor %]
\end{longtable}

[% if checklists %]
% --- Checklist appendix ---
\appendix
[% if taal == 'en' %]
\section{Checklist Results}
[% else %]
\section{Checklist resultaten}
[% endif %]

[% for cl in checklists %]
\subsection{[[ cl['naam'] | latex_escape ]] ([[ cl['type'] | latex_escape ]]) --- [[ cl['target'] | latex_escape ]]}

\noindent\textbf{[[ cl['stats']['pass'] ]] pass, [[ cl['stats']['fail'] ]] fail, [[ cl['stats']['warn'] ]] warn, [[ cl['stats']['skip'] ]] skip, [[ cl['stats']['na'] ]] n/a, [[ cl['stats']['open'] ]] open}
\vspace{0.5em}

[% for phase in cl['phases'] %]
\subsubsection{[[ phase['title'] | latex_escape ]]}

\begin{longtable}{|p{2.5cm}|p{9cm}|p{3cm}|}
  \hline
  \textbf{Ref} & \textbf{Item} & \textbf{Status} \\
  \hline
  \endhead
[% for item in phase['items'] %]
  \texttt{\small [[ item['ref'] | latex_escape ]]} & [[ item['title'] | latex_escape ]] & [% if item['status'] == 'pass' %]\textcolor{ogreen}{PASS}[% elif item['status'] == 'fail' %]\textcolor{ored}{FAIL}[% elif item['status'] == 'warn' %]\textcolor{oorange}{WARN}[% elif item['status'] == 'skip' %]SKIP[% elif item['status'] == 'na' %]N/A[% else %]OPEN[% endif %] \\
  \hline
[% endfor %]
\end{longtable}

[% endfor %]
[% endfor %]
[% endif %]

\end{document}
